% Study sheet: Volume II, Chapter 10, Serving at Scale.
% Numbers: \val{...} from inference_studysheet_numbers.tex, written by inference_studysheet_numbers.py.
\documentclass{studysheet}
\input{inference_studysheet_numbers}
\sheetvolume{2}
\sheetchapter{10}{Serving at Scale}
\sheetsources{Numbers from the mlsysim registry and the chapter's modeled scenarios; HBM in GiB, other sizes decimal; assumed inputs are marked.}

\begin{document}
% ===================================================================== FRONT
\frontside
\begin{bigidea}
Training is paid once; serving, on \hl{every request}. A fleet's
capacity is its \hl{SLO goodput}: completed requests that meet their first-token (TTFT) and per-token (TPOT) targets.
In C\textsuperscript{3} terms, batching's compute reuse must pay for the communication and coordination
it adds, inside the latency budget.
\end{bigidea}

\begin{leftcol}
\sheetsec{Key ideas}
\begin{ideas}
\idea{Serving cost recurs.}{At high traffic it overtakes training within weeks.}
\idea{A late answer counts for nothing.}{Network, serialization, coordination and queueing
  share the budget model execution leaves.}
\idea{Batching trades waiting for reuse.}{Shared weight reads lift throughput toward a ceiling; a
  batching window caps formation delay, not queueing.}
\idea{The KV cache is the admission budget.}{Each request's cache grows with context;
  continuous batching and paging use memory better, but never shrink the cache.}
\idea{Prefill computes; decode streams weights.}{Chunked prefill feeds a long prompt one chunk per
  iteration: token gaps shrink; the prompt's first token waits.}
\idea{Distribution sends a bill.}{Per request, the Fleet Law holds: sharding adds collectives, so
  it stays in one NVLink domain. Replicas, not shards, add capacity.}
\end{ideas}

\vspace{10pt}
\sheetsec{Which limit binds? Read the evidence}
\begin{sheettable}{@{}l L L@{}}
\thead{Limit} & \thead{Evidence} & \thead{Lever}\\ \midrule
Weight fit & weights exceed one device's HBM & lower-bit weights; shard inside the node\\
KV capacity & admission stalls while compute idles & fewer KV heads, 8-bit cache, shorter context\\
Prefill interference & token gaps spike as long prompts arrive & chunked prefill, or separate pools\\
Queueing & p99 grows with load, model time steady & add replicas; admit less while they start\\
\end{sheettable}

\vspace{10pt}
\fcard{Lifetime cost}%
  {\ModelLifetimeCost = \TrainingCost + \ServingCostPerQuery \times \DeploymentDuration \times \ArrivalRate}%
  {Training, plus cost per request $\times$ time $\times$ request rate. \val{users}M users
   $\times$ \val{perday} requests a day at \$\val{perrequestcost}: \$\val{servingm}M a year, \val{ratio}$\times$ the \$\val{trainm}M training.}%
  {}{}
\end{leftcol}%
\begin{rightcol}
\sheetsec{Formulas you will use}
\fcard{Batching curve}%
  {\BatchServiceTime(\BatchSize) = \BatchFixedOverheadTime + \BatchSize\,\MarginalRequestServiceTime,
   \qquad \EffectiveServiceRate = \frac{\BatchSize}{\BatchServiceTime(\BatchSize)}}%
  {Throughput rises toward $1/\MarginalRequestServiceTime$; at the knee $\BatchSize \approx
   \BatchFixedOverheadTime/\MarginalRequestServiceTime$ it reaches half. Decode: $\BatchFixedOverheadTime$ = \val{tfixed} ms
   streams the weights, $\MarginalRequestServiceTime$ = \val{treq} ms reads one sequence's cache.}%
  {A decode step yields one token per stream: $\BatchServiceTime$ must meet the per-token target.}%
  {\begin{tikzpicture}[sheet]
     \draw[dim] (20,0) -- (226,0);
     \draw[dim] (20,0) -- (20,78);
     \foreach \b in {0,80,200,400}{
       \draw[dim] ({20+0.5*\b},0) -- ({20+0.5*\b},-2.5);
       \node[small, anchor=north] at ({20+0.5*\b},-4) {\b};}
     \node[small, anchor=north] at (150,-13) {batch size (sequences)};
     \node[small, anchor=east] at (17,35) {{1,000}};
     \node[small, anchor=east] at (17,70) {{2,000}};
     \node[small, anchor=south west] at (0,80) {tokens/s};
     \draw[accent, dashed, line width=0.6pt] (20,70) -- (220,70);
     \node[note, anchor=south east] at (220,71) {ceiling $1/\MarginalRequestServiceTime$};
     \edef\bplot{\noexpand\draw[accent, line width=1pt] plot coordinates {\val{curve}};}\bplot
     \draw[dim, dashed] (\val{xknee},0) -- (\val{xknee},\val{yknee});
     \filldraw[fill=accent, draw=accent] (\val{xknee},\val{yknee}) circle (2pt);
     \node[lab, anchor=north west] at (\val{xknee}+5,\val{yknee}-2) {knee \val{knee}: half the ceiling};
   \end{tikzpicture}}
\fcard{KV admission budget}%
  {\begin{aligned}\KvCacheBytesPerToken &= 2\,\LayerCount\,\KvHeadCount\,\AttentionHeadDimension\,\ElementStorageSize\\[2pt]
   \BatchSize_{\max} &= \left\lfloor \frac{\text{HBM} - \WeightMemory - \text{reserve}}{\SequenceLength\,\KvCacheBytesPerToken} \right\rfloor,
   \quad \WeightMemory = \Parameters\,\frac{\NumericBitWidth}{8}\end{aligned}}%
  {Key and value, $\LayerCount$ layers, $\KvHeadCount$ KV heads of size $\AttentionHeadDimension$, $\ElementStorageSize$ bytes each;
   $\SequenceLength$ context tokens; $\Parameters$ parameters at $\NumericBitWidth$ bits.}%
  {An H100 holds \val{hbmgib} GiB = \val{hbmgb} GB.}{}
\fcard{Fan-out tail}%
  {\Pr(\text{request slow}) = 1 - (1 - \IndependentEventProbability)^{\SeriesComponentCount}}%
  {A request waits for its slowest of $\SeriesComponentCount$ leaves, each slow with probability
   $\IndependentEventProbability$. \val{leaves} leaves at p\val{leafpct}: \val{slowpct}\% of requests slow;
   holding p\val{leafpct} needs each leaf slow under \val{leafbudget}\%.}%
  {Levers: fewer shards; route on queue length or hedge (lower $\IndependentEventProbability$); a deadline with a partial result.}{}
\end{rightcol}

% ===================================================================== BACK
\backside
\begin{answers}
\ans{Finish}{INT8 cache: \val{pertoken8} B/token, \val{perrequest8} GB per request; \val{budget} $\div$ \val{perrequest8}
  $\to$ \res{\val{batchkveight}}. INT4 weights: \val{weightsfour} GB, so \val{nodegb} $-$ \val{weightsfour} $-$ \val{reservegb}
  = \val{budgetfour} GB; $\div$ \val{perrequest} $\to$ \res{\val{batchweightsfour}}. The cache wins: halving per-request
  state doubles the count; INT4 frees only the weight share.}
\ans{1}{Knee \val{qfixed} $\div$ \val{qreq} = \val{qknee}; ceiling one over \val{qreq} ms = \val{qceiling} tokens/s.
  At \val{qbatch}: step \val{qfixed} + \val{qextra} = \val{qstep} ms, \val{qrate} tokens/s (\val{qshare}\% of ceiling),
  but \val{qstep} ms misses \val{qtarget} ms. Largest batch: (\val{qtarget} $-$ \val{qfixed}) $\div$ \val{qreq} $\to$
  \val{qmax} (\val{qmaxstep} ms). The target, not the knee, sets the batch (batching card).}
\columnbreak
\ans{2}{One minus \val{qfast} to the power \val{qleaves} = \val{qslowpct}\% slow. Leaf budget: one minus \val{qfast} to the power one over \val{qleaves} =
  \val{qleafbudget}\%: p\val{qleafpct}. Levers: fewer, larger shards; route on queue length; hedge slow leaves;
  deadline with a partial result (fan-out card).}
\ans{3}{A monolithic prefill lands in every stream's token gap. \val{prompt} $\div$ \val{chunk}
  = \val{chunkratio}, so \val{chunkiters} iterations: each gap is one hybrid step, but the prompt's first token waits all
  \val{chunkiters}.}
\ans{4}{\val{qusers}M $\times$ \val{qper} = \val{qdaily}M a day $\times$ \val{qdays} = \val{qannual}B requests $\times$
  \$\val{qcost} = \$\val{qservem}M, \val{qratio}$\times$ training. At \$\val{qdaycost} a day, serving matches training
  after \val{qcross} days; a \val{qcut}\% cut saves \$\val{qsave}M a year (cost card).}
\end{answers}

\begin{leftcol}
\sheetsec{Worked example}
\begin{practicebox}\boxtitle{How many \val{contextk}K-token requests fit on one \val{gpus}-GPU node?}
A \val{params}B-parameter model: \val{layers} layers, \val{kvheads} KV heads of dimension \val{headdim}, FP16
(\val{fpbytes} bytes) weights and cache. \val{gpus} H100s $\times$ \val{hbmgib} GiB = \val{nodegib} GiB
(\val{nodegb} GB); reserve \val{reservegb} GB for runtime.
\begin{steps}
\item Weights: \val{params}B $\times$ \val{fpbytes} B = \res{\val{weightsgb} GB}.
\item Per token: $2 \times \val{layers} \times \val{kvheads} \times \val{headdim} \times \val{fpbytes}$ = \val{pertoken} B;
  $\times$ \val{context} tokens = \res{\val{perrequest} GB} per request.
\item Free for cache: \val{nodegb} $-$ \val{weightsgb} $-$ \val{reservegb} = \res{\val{budget} GB}.
\item $\BatchSize_{\max} = \lfloor \val{budget} \div \val{perrequest} \rfloor$ = \res{\val{maxbatch}}.
\end{steps}\vspace{4pt}
{\centering
\begin{tikzpicture}[sheet]
  \pgfmathsetmacro\a{\val{barweights}} \pgfmathsetmacro\b{\a+\val{barreserve}} \pgfmathsetmacro\c{\b+\val{barkv}}
  \draw[accbox] (0,0) rectangle (\a,14); \node[lab] at ({\a/2},7) {weights};
  \draw[nbox] (\a,0) rectangle (\b,14);
  \draw[hotbox] (\b,0) rectangle (\c,14);
  \foreach \i in {1,...,11}{\draw[white, line width=0.4pt] ({\b+\i*\val{barone}},0) -- ({\b+\i*\val{barone}},14);}
  \draw[nbox] (\c,0) rectangle (230,14);
  \node[note, anchor=north] at ({(\b+\c)/2},-3) {\val{maxbatch} requests $\times$ \val{perrequest} GB};
  \node[note, anchor=north west] at (0,-3) {\val{weightsgb} GB};
  \node[note, anchor=south east] at (230,16) {node: \val{nodegb} GB};
  \node[note, anchor=south west] at (0,16) {white: reserve, unused};
\end{tikzpicture}\par}
\takeaway{The cache, not the weights, sets the batch.}
\end{practicebox}

\vspace{10pt}
\sheetsec{Your turn: finish it}
\begin{practicebox}\boxtitle{Compress the cache or the weights?}
Same node and reserve. A: INT8 cache (one byte). B: INT4 weights (half a byte per parameter).\par\vspace{4pt}
\textbf{Given:} A halves $\ElementStorageSize$; B shrinks only $\WeightMemory$.\par
\textbf{You:} find each batch, and why one lever admits more.\par\vspace{24pt}
\end{practicebox}
\end{leftcol}%
\begin{rightcol}
\sheetsec{Common mistakes}
\mistake{``The model fits on the GPU, so it can serve long contexts.''}%
  {INT4 weights (\val{weightsfour} GB) fit one \val{hbmgb} GB H100, and one \val{contextk}K request brings it to
   \val{onefit} GB; a second makes \val{twofit} GB, past \val{hbmgb}.}{}
\mistake{``Rising p99 means we need a bigger, sharded replica.''}%
  {If model time is flat while queueing grows, demand has outgrown the replicas. Sharding only divides each
   request's compute and adds collectives: add replicas, and cap admission during each one's \val{coldmin}-min cold start.}{}
\mistake{``Size the fleet for average throughput.''}%
  {At \val{util}\% utilization, \val{growth}\% more traffic gives \val{util} $\times$ \val{growfactor} = \val{newutil}\%:
   no headroom for a burst. Size for peaks.}{}

\vspace{2pt}
\sheetsec{Check yourself}
\begin{qs}
\item \assumed{} a decode step costs \val{qfixed} ms plus \val{qreq} ms per sequence; per-token target \val{qtarget} ms.
  Find the knee and ceiling. Batch \val{qbatch}: what share of the ceiling, and does it meet the target? What batch does?
\item \assumed{} a request fans out to \val{qleaves} leaves, each at its own p\val{leafpct}. What share of requests is
  slow? What leaf percentile holds p\val{leafpct}, and which levers get there?
\item Average per-token time is fine, but token gaps spike when long prompts arrive. Why? What do \val{chunk}-token
  chunks buy and cost for a \val{prompt}-token prompt?
\item \assumed{} \$\val{qtrainm}M training; \val{qusers}M users send \val{qper} requests a day at \$\val{qcost}. First-year
  serving cost and ratio? When does serving match training? What does a \val{qcut}\% unit-cost cut save?
\end{qs}
\end{rightcol}
\end{document}
